\section{语用分析}
\label{sec:nlp-analysis-pragmatics}

“明天能完成吗”在制定计划时是在询问可行性；对方已经两次延误之后，同一句话更像催促。字面命题几乎相同，在交流中起的作用却不同。\term{语用分析}（pragmatic analysis）研究话语在具体背景下的作用：它在交流中完成了什么行为，省略了哪些可以从背景补出的内容，以及说话者通过字面意义暗示了什么。这里的背景包括说话者与听话者、当前话语回应的对象、此前的对话历史，以及双方共享的常识。

本节的三个小节对应三类问题。话语行为把每一轮话语归入请求、告知、确认等功能类别，是对话系统理解用户与规划回复的基础；省略恢复与话语改写把依赖历史才能理解的话语补全为独立完整的句子，使下游的检索与问答模块可以只看当前句；会话含义与间接言语行为则处理字面意义与说话者意图不一致的情形，这是语用中最依赖推理的部分。三者的共同难点是，所需的信息不在当前句之内，模型必须从背景中取用，又不能补出背景没有支持的内容。

\subsection{话语行为}
\label{sec:nlp-analysis-dialogue-acts}

\term{话语行为}（dialogue act）用请求、告知、确认、承诺等标签表示一次话语在交流中的功能。这一概念来自言语行为理论（见第~\ref{sec:nlp-analysis-indirect-acts}小节）：说话不只是陈述命题，也是在做事，如提问、许诺、请求。话语行为标注把这种思想落实为可操作的标签集。Switchboard对话行为语料（SwDA）在1155段电话交谈上按SWBD-DAMSL规范标注，常用的归并标签集有42类，包括陈述观点、陈述非观点、附和（backchannel，如“嗯”“对”）、是非问、同意等\cite{nlp-analysis-jurafsky1997swbd}；ICSI会议对话行为语料（MRDA）在75场多人会议上标注，原始标签由11个一般标签和39个具体标签组合而成\scite{nlp-analysis-shriberg2004}；后续研究最常用的是随语料发布的五类映射：陈述、提问、话轮控制（floor mechanism，如抢话与占话）、附和与打断。两份语料的单位都是切分后的话语片段，而不是完整话轮：一个话轮可以包含“嗯，对，我觉得周五不行”这样连续的几个片段。多功能问题在标注规范层面也有专门的处理：国际标准ISO 24617-2把话语功能分为任务、反馈、话轮管理、时间管理等多个维度，一个片段可以在每个维度上各有一个功能，例如“嗯，周五不行”同时承担反馈维度的确认收听与任务维度的告知\scite{nlp-analysis-bunt2012iso}。采用这类规范时，识别就成为每个维度各一个分类器的多任务问题。

话语行为识别有四个挑战。同一句式在不同前文下功能不同，“明天能完成吗”既可以是询问也可以是催促。说话者转换：对话行为的依赖往往跨越说话者，问句之后通常是回答，而回答出自另一方。多功能话语：“谢谢，能再发一份吗”同时致谢与请求，标签集是否允许多标签要事先规定。标签集差异：SwDA与MRDA的类别粒度、单位切分都不同，模型和结论不能直接迁移。此外，在线识别只能使用当前及以前的话轮，事后分析可以看到整段对话，两种设置的输入条件不同，结果不可直接比较。

方法谱系按怎样使用上下文展开。单句分类只看当前话语，用词语、句式、说话者等特征送入最大熵、SVM或卷积网络分类器，这是最直接的基线，但它对上面第一项挑战无能为力。跨轮HMM与CRF把行为序列看作马尔可夫链，用行为之间的转移刻画前后依赖。层次对话编码先在话轮内汇聚词语得到话轮向量，再在话轮之间编码对话历史，用学到的表示代替人工特征，并保留CRF对标签依赖的建模。

单句分类器也可以用来扩展标签集。新领域出现训练数据中没有的行为类别时，可以把每个标签的定义与几个例子写成一段描述，用句对编码器计算“当前话语--标签描述”的匹配分数，在所有描述中取最高者。这改变的是类别信息的来源：模型不再从大量样本中学习一个类别，而是从描述中读出它的边界，因此描述必须写清相近类别的区别，例如“询问可行性”与“要求加快进度”。训练与预测时可见的信息也应一致：若标注者看了后续回答才判断当前行为，而系统在预测时只能看到当前及以前的话语，监督信号就含有预测时不可得的信息，这种任务应明确为事后分析。

\subsubsection{跨轮HMM与CRF话语行为标注}
\label{sec:nlp-analysis-dialogue-sequence}

Stolcke等把一段对话的行为序列建模为HMM：每轮行为$a_t$是隐状态，话语$u_t$是观测，联合概率为$\prod_t p(a_t\mid a_{t-1})\,p(u_t\mid a_t)$\scite{stolcke-etal-2000-dialogue}。转移概率由已标注的相邻行为计数估计，相当于一个行为层面的$n$元语言模型；发射概率由每类行为各自的词语$n$元语言模型给出，还可以结合韵律特征。若手头只有判别式分类器给出的$p(a_t\mid u_t)$，可以按贝叶斯公式除以先验$p(a_t)$，换算为与$p(u_t\mid a_t)$成比例的量后再代入，不能直接把后验当作发射概率使用。

固定观测“明天能完成吗”，设教学发射概率为$p(u\mid I)=0.6$、$p(u\mid U)=0.4$，I为询问，U为催促。前轮行为是计划讨论P时，转移到I、U的概率为$0.8,0.2$，乘积为$0.48,0.08$，归一化后I占$6/7$；前轮是延误解释E时，转移概率为$0.2,0.8$，乘积为$0.12,0.32$，U占$8/11$。观测的贡献没有变，改变排序的是行为上下文。

前轮行为本身也不确定时，在线识别用前向算法逐轮更新过滤分布：
\[
  p(a_t\mid u_{\leq t})\propto p(u_t\mid a_t)\sum_{a_{t-1}}p(a_t\mid a_{t-1})\,p(a_{t-1}\mid u_{<t}).
\]
事后分析整段对话时，前向--后向算法给出每轮的边缘分布，Viterbi算法给出整条最可能的行为序列。两者回答不同的问题：前者对前一轮的所有可能行为求和，后者只沿最好的一条路径取最大，因此Viterbi路径上某一轮的行为，不一定是该轮边缘概率最大的行为。看一个两轮的例子，第二轮仍是“明天能完成吗”，沿用上面的转移与发射概率。若第一轮的过滤分布为P占$0.55$、E占$0.45$，第二轮I的未归一化概率为$0.55\times0.48+0.45\times0.12=0.318$，U为$0.55\times0.08+0.45\times0.32=0.188$，边缘偏向I；Viterbi中到达I的最优路径来自P，分数$0.264$，到达U的来自E，分数$0.144$，最优路径P--I，两者一致。若第一轮改为P占$0.35$、E占$0.65$，第二轮I的未归一化概率为$0.35\times0.48+0.65\times0.12=0.246$，U为$0.35\times0.08+0.65\times0.32=0.236$，边缘仍略偏向I（约$0.51$）；Viterbi比较的却是$0.35\times0.48=0.168$与$0.65\times0.32=0.208$，最优路径为E--U。I的概率来自P与E两条前驱路径之和，单看任何一条都不如E--U高。需要给每一轮输出最可能的标签时用边缘分布，需要一条前后一致的行为序列时用Viterbi。HMM的定义与前向后向、Viterbi递推见第~\ref{chap:nlp-classical-models}章。

线性链CRF直接对条件概率建模，可以把话语文字、说话者是否转换、与前一轮的词语重叠等任意特征放进位置得分，并学习行为之间的转移分数，其定义、梯度与解码见第~\ref{sec:nlp-classical-crf}小节。对同一问句设I、U的位置分数为$1,0.5$；P之后的转移分数为$1,0$，两候选总分为$2,0.5$，偏向I；E之后的转移分数为$-0.5,1$，总分为$0.5,1.5$，偏向U。与分词的BMES标注不同，话语行为之间一般没有必须禁止的转移，因此不需要第~\ref{sec:nlp-analysis-paradigm-tagging}小节那样的合法转移掩码，转移矩阵完全由数据学习。

HMM与CRF的局限在于，上下文只通过行为标签传递：两段历史若行为标签序列相同而内容不同，模型无法区分。“明天能完成吗”之前若是“上次你说周三能好”，催促的意味来自具体内容，而这一内容在“告知”这个标签中已经丢失。

\subsubsection{层次对话编码与分层BiLSTM-CRF}
\label{sec:nlp-analysis-dialogue-hierarchical}

层次对话编码让历史话语的内容而不仅是标签进入当前判断。对话天然有两层结构：词组成话语，话语组成对话。Kumar等的分层BiLSTM-CRF按这两层各用一个双向LSTM编码，再接线性链CRF解码整段行为序列\cite{nlp-analysis-kumar2018}，结构如图~\ref{fig:nlp-analysis-hierarchical-dialogue-act}所示。

设一段对话$C$含$R$个话语$u_1,\ldots,u_R$，第$j$个话语有$S_j$个词。第一层是词级编码：每个词经嵌入层得到$\bm e_k\in\mathbb R^{d_e}$，词级BiLSTM输出每个位置的前后向隐状态拼接$\bm h_k\in\mathbb R^{2h_1}$。话语向量取最后一个时间步的表示，$\bm v_j=\bm h_{S_j}\in\mathbb R^{2h_1}$；原文也比较了平均池化，最终采用末状态。第二层是话轮级编码：把$\bm v_1,\ldots,\bm v_R$当作一个序列送入另一个BiLSTM，输出$\bm g_j\in\mathbb R^{2h_2}$，它同时包含$u_j$之前与之后的话语内容。第三层是CRF：设行为标签集大小为$K$，位置分数$\bm o_j=\bm W\bm g_j+\bm b\in\mathbb R^{K}$，转移矩阵$\bm A\in\mathbb R^{K\times K}$与位置无关，行为序列$\bm y$的条件概率为
\begin{equation}
  p(\bm y\mid C)=\frac{\exp\bigl(\sum_{j=1}^{R}(\bm o_{j,y_j}+\bm A_{y_{j-1},y_j})\bigr)}
   {\sum_{\bm y'}\exp\bigl(\sum_{j=1}^{R}(\bm o_{j,y'_j}+\bm A_{y'_{j-1},y'_j})\bigr)}.
  \label{eq:nlp-analysis-hierarchical-crf}
\end{equation}
训练最大化各段对话的对数似然$\sum_C\log p(\bm y\mid C)$，嵌入层、两层BiLSTM与CRF参数端到端联合更新；推断用Viterbi算法求最优行为序列。用两轮对话看CRF层的作用（数值为教学构造）。设标签只有问（Q）与答（A）两类，两轮话语的位置分数为$\bm o_1=(2,0)$、$\bm o_2=(0.4,0.6)$，即第二轮单看自身略偏向答；转移矩阵中$\bm A_{Q,A}=1$、$\bm A_{Q,Q}=-1$、$\bm A_{A,Q}=\bm A_{A,A}=0$，起始转移取0。四条序列的总分为QQ：$2+0.4-1=1.4$，QA：$2+0.6+1=3.6$，AQ：$0+0.4=0.4$，AA：$0+0.6=0.6$。最优序列QA的概率为$e^{3.6}/(e^{1.4}+e^{3.6}+e^{0.4}+e^{0.6})\approx0.83$。若去掉转移矩阵、逐轮独立取最大，第二轮也会选A，但概率只有约$0.55$；转移矩阵把“问之后多为答”的规律加了进来，使判断更确定。若第二轮的位置分数明显偏向Q，例如$(2,0)$，QQ得$3$，QA得$3$，两者持平，此时是否连续两问取决于话语内容与转移规律的权衡。按张量形状看，一个批次中的对话先被展开为话语的批次，形状为（话语数，最大词数，$d_e$），词级编码后得到（话语数，$2h_1$），再按对话重新组织为（对话数，最大话语数，$2h_1$）送入话轮级编码。

\begin{figure*}[htbp]
  \centering
  % 分层BiLSTM-CRF话语行为标注：词级编码汇聚为话轮向量，话轮级编码建模轮间依赖，CRF解码行为序列。
\begin{tikzpicture}[
  x=1mm,y=1mm,figure picture,
  cell/.style={figure node,minimum width=12mm,minimum height=6mm,
    inner sep=1pt,font=\figurefont\scriptsize},
  word/.style={font=\figurefont\scriptsize},
  row label/.style={font=\figurefont\scriptsize,text=FigureMuted,anchor=east},
  flow/.style={figure flow},
  link/.style={draw=FigureStroke,line width=.7pt,{Stealth[length=1.3mm]}-{Stealth[length=1.3mm]}},
  pair/.style={draw=FigureOutput,line width=.9pt}
]
  \foreach \j/\x/\txt in {1/0/{计划定在周五},2/36/{明天能完成吗},3/72/{可以}}{
    \node[word] (u\j) at (\x,0) {$u_{\j}$：\txt};
    \node[cell,draw=FigureInput,minimum width=26mm] (w\j) at (\x,10)
      {词级BiLSTM，取末状态};
    \node[cell,draw=FigureInput] (v\j) at (\x,21) {$\bm v_{\j}$};
    \node[cell,draw=FigureModel] (g\j) at (\x,33) {$\bm g_{\j}$};
    \node[cell,draw=FigureOutput] (y\j) at (\x,46) {$y_{\j}$};
    \draw[flow] (u\j) -- (w\j);
    \draw[flow] (w\j) -- (v\j);
    \draw[flow] (v\j) -- (g\j);
    \draw[flow] (g\j) -- (y\j);
  }
  \draw[link] (g1) -- (g2);
  \draw[link] (g2) -- (g3);
  \draw[pair] (y1) -- node[above,font=\figurefont\scriptsize] {$\bm A_{y_1,y_2}$} (y2);
  \draw[pair] (y2) -- node[above,font=\figurefont\scriptsize] {$\bm A_{y_2,y_3}$} (y3);
  \node[row label] at (-16,21) {话轮向量};
  \node[row label] at (-16,33) {话轮级BiLSTM};
  \node[row label] at (-16,46) {线性链CRF};
\end{tikzpicture}

  \caption{分层BiLSTM-CRF话语行为标注。每个话语先经词级BiLSTM编码，取末状态作为话语向量$\bm v_j$；话轮级BiLSTM在话语向量序列上编码得到$\bm g_j$，使每个位置都含有前后话语的内容；线性链CRF以$\bm g_j$给出位置分数，以转移矩阵$\bm A$给出相邻行为的分数，联合解码整段行为序列。}
  \label{fig:nlp-analysis-hierarchical-dialogue-act}
\end{figure*}

两层编码各自承担一项挑战。词级编码把话语内容压缩为向量，解决了HMM与CRF只能传递标签的问题；话轮级编码让前后话语的内容互相影响，CRF的转移矩阵再刻画“问句之后多为回答”这类标签依赖。原论文在SwDA（42类）与MRDA（5类）上评测，准确率分别为79.2\%与90.9\%，均高于此前报告的结果；训练集分别有1003段与51段对话。两份语料的标签分布都极不均衡，SwDA中陈述非观点与附和两类就占一半以上，MRDA中陈述类也超过一半，因此只看准确率容易高估模型对少数类的识别能力，宜同时报告宏平均F1。

这一模型有两点需要按任务调整。其一，话轮级BiLSTM的后向部分读取了未来话语，只适用于事后分析；在线识别必须改用单向编码，或在每一轮只对已发生的前缀重新编码。其二，原模型没有区分说话者，而说话者是否转换对行为预测很有信息量（同一人连续两句多为陈述的延续，换人之后常是回答或附和）。常见的扩展是在$\bm v_j$上拼接说话者嵌入，或拼接“与上一话语是否同一说话者”的指示特征；在多方对话中，当前话语回应的未必是紧邻的上一句，还需要输入或同时预测回应关系。多功能话语若规范允许多个标签，应把CRF换成逐标签的sigmoid输出，或先切分片段再分别标注。

后续工作主要替换编码器。Raheja与Tetreault在话语层与对话层引入自注意力，让话语向量按注意力汇聚词语表示，而不是只取末状态，并保留CRF解码\cite{nlp-analysis-raheja2019}；用预训练编码器替换词级BiLSTM时，把当前话语与前几轮拼接后编码，可以同时覆盖词级与话轮级两层。这些替换改变的是表示能力，话语—对话的两层结构与CRF的标签依赖仍然沿用。

表~\ref{tab:nlp-analysis-dialogue-act-methods}对照三类方法。

\begin{table*}[htbp]
  \centering
  \small
  \begin{tabularx}{\linewidth}{@{}l XXX@{}}
    \toprule
    方法 & 解决的问题 & 核心机制 & 新的代价或局限\\
    \midrule
    单句分类 & 话语功能的基本识别 & 词语、句式、说话者特征加分类器 & 不用上下文，同一句式只能有一个答案\\
    HMM与CRF序列标注 & 前后话语行为的依赖 & 行为转移概率或转移分数，前向或Viterbi解码 & 上下文只经标签传递，丢失历史内容\\
    分层BiLSTM-CRF & 历史内容进入当前判断 & 词级与话轮级两层编码，CRF建模标签依赖 & 双向编码只适合事后分析；需另加说话者信息\\
    \bottomrule
  \end{tabularx}
  \caption{话语行为识别方法对照。}
  \label{tab:nlp-analysis-dialogue-act-methods}
\end{table*}

\subsection{省略恢复与话语改写}
\label{sec:nlp-analysis-utterance-rewriting}

多轮对话中，用户经常省略前文已经说过的内容。历史H1为“今天能到图书馆借书吗”，历史H2为“今天的空气质量如何”，当前话语都是“那明天呢”。在H1之下，它的意思是“明天能到图书馆借书吗”；在H2之下，是“明天的空气质量如何”。检索系统、问答模型若只拿到“那明天呢”，几乎无法工作。\term{话语改写}（utterance rewriting）把当前话语改写为脱离历史也能理解的完整句子，补出被省略的成分（\term{省略恢复}，ellipsis recovery），并把代词替换为它所指的内容。改写之后，下游模块可以把对话问题当作单轮问题处理。

这一任务有几份常用数据。CANARD由Elgohary等构建：从对话式问答数据QuAC中取出依赖上下文的问题，请标注者把每个问题改写为独立完整的问题，共约4万条\scite{nlp-analysis-elgohary2019}。TREC CAsT（对话式辅助检索评测）面向对话式段落检索，每个话题是一串相关的问题，后面的问题大量省略和指代前面的内容，评测方同时提供人工改写的独立问题，便于比较“用原始问题检索”与“用改写问题检索”的差距\scite{nlp-analysis-dalton2020}。中文方面，Pan等构建了Restoration-200K，在约20万条开放领域多轮闲聊上标注了补全后的话语\cite{nlp-analysis-pan2019restoration}；Su等收集了2万段中文多轮对话并为其中的话语标注改写，他们对对话数据的分析发现，约七成的话语存在指代或省略\scite{nlp-analysis-su2019rewriter}。

省略恢复的挑战有三项。其一，省略的位置不显式：当前话语中没有占位符告诉模型缺了什么、缺在哪里。其二，补全内容必须来自历史：“明天能到图书馆借书吗”中的“图书馆”“借书”都应能在H1中找到来源，模型补出历史中没有的内容，就改变了用户的问题。其三，历史不足时应保留未知：若对话没有历史，“那明天呢”就没有可靠的补全，应当保持原样或请求澄清。评价也因此有两个层面：改写与参考改写的重合度（BLEU、ROUGE、精确匹配，或只看补入词的精确率与召回率），以及改写对下游检索或问答效果的提升。重合度指标中，只看补入词的精确率与召回率最能反映省略恢复本身：设参考改写相对当前话语新增的词集合为$A$，系统改写新增的词集合为$B$，则精确率为$|A\cap B|/|B|$，召回率为$|A\cap B|/|A|$。“那明天呢”在H1下的参考改写新增“能、到、图书馆、借书、吗”，系统若改写为“明天图书馆开门吗”，新增“图书馆、开门、吗”，精确率为$2/3$，召回率为$2/5$。BLEU在这里不够敏感：改写与当前话语大量重叠，照抄当前话语也能拿到不低的分数。下游评价则直接用改写去检索，在TREC CAsT上常报告前3个结果的NDCG@3，并与人工改写的检索效果比较。

方法谱系从结构化场景走向开放领域。改写并不是利用历史的唯一方式。也可以把完整历史与当前话语一起交给下游模型，让它在内部完成补全。显式改写的好处是中间结果可读、可检查，可以接入只接受单轮输入的现成检索系统，改写模型与下游模块可以分别优化；代价是多了一个模块，改写的错误会传给下游，且改写可能丢掉历史中对下游有用、却不属于省略内容的信息。两种方式常结合使用：检索用改写后的问题，生成回答时仍提供完整历史。

候选承接与槽位继承依赖预定义的槽位结构，把上一轮的槽值按规则或分类器继承到当前轮，可靠但覆盖有限；复制式改写用序列到序列模型从历史和当前句两路复制词语，不需要槽位定义，覆盖开放领域；大语言模型改写用提示直接完成，不需要专门的训练数据，但可能补入历史中没有的信息，需要来源检查。

\subsubsection{候选承接与槽位继承}
\label{sec:nlp-analysis-slot-carryover}

在任务型对话中，用户的意图由一组槽位描述，例如查询活动安排时有“活动”“日期”“时段”“时间”“地点”等槽位。每一轮结束后，系统维护一个对话状态，即目前已确定的槽值集合。省略在这种场景下表现为当前话语只给出部分槽值，其余槽值需要从对话状态继承。\term{槽位继承}（slot carryover）把这一判断写成分类问题：对历史中出现过的每个候选槽值$(s,v)$，判断它在当前轮是否仍然有效\scite{nlp-analysis-naik2018carryover}。

设当前话语为$u_t$，候选槽值$(s,v)$最近一次出现在$\delta$轮之前。编码器分别给出当前话语的表示$\bm q_t$、候选槽名与槽值的表示$\bm c_{s,v}$，再加上距离嵌入$\phi(\delta)$，以sigmoid输出继承概率：
\[
  p(\mathrm{carry}=1\mid u_t,s,v)=\sigma\bigl(\bm w^{\top}\operatorname{FFNN}([\bm q_t;\bm c_{s,v};\bm q_t\circ\bm c_{s,v};\phi(\delta)])\bigr),
\]
以标注的继承标签做二元交叉熵训练。例如当前话语“那上午呢”与候选｛时段：下午｝的匹配表示显示两者属于同一槽且取值冲突，网络会给出很低的继承概率；对候选｛活动：读书会｝，当前话语没有提及任何活动，距离$\delta=1$，继承概率高。当前话语若显式给出了某个槽的新值，该槽的旧值直接被替换，不再参与判断。规则方法是这一模型的特例：“当前轮未提及的槽一律继承”相当于对所有候选输出1，它在简单场景下够用，但无法处理槽之间的依赖。

依赖关系最常出现在时间类槽位上。设读书会的通知为“读书会改到周五下午两点”，对话状态为｛活动：读书会，日期：周五，时段：下午，时间：14:00｝。用户接着问“那地点呢”：只是换了询问对象，四个槽全部继承。用户问“改到周六呢”：日期被替换为周六，而“下午两点”是否随之保留，取决于用户是在问“周六同一时间”还是在问“周六有没有安排”，分类器应给出中等的继承概率，系统可以确认。用户问“那上午呢”：日期继承为周五，时段被替换为上午，精确时间14:00与上午矛盾，必须丢弃而不能继承。这说明继承判断不能逐槽独立进行，替换一个粗粒度的槽（时段）会使依附于它的细粒度槽（时间）失效。

槽位继承的长处是输出结构明确、可以逐槽检查，代价是只能处理槽位体系覆盖的内容。“那明天呢”在借书场景下需要补出“到图书馆借书”这一动作，在开放领域对话中省略的可能是任意短语，没有预先定义的槽位可以承接。

\subsubsection{基于复制的话语改写}
\label{sec:nlp-analysis-copy-rewriting}

开放领域的改写需要从历史中取出任意片段。观察可知，改写结果中的词几乎都来自两个地方：要么来自当前话语本身（“那明天呢”中的“明天”），要么来自历史（“图书馆”“借书”）。Su等据此设计了只做复制、不从词表生成的改写模型：解码器每一步从历史或当前话语中各自复制一个词，再由门控决定取哪一路\scite{nlp-analysis-su2019rewriter}。复制机制的一般形式见式~\eqref{eq:nlp-analysis-copy}，这里是它的双源版本。

编码器是Transformer，历史H与当前话语U拼接后共同编码，每个词元除词嵌入与位置嵌入外，还加上一个轮次嵌入表示它属于哪一轮，得到编码$\bm E\in\mathbb R^{(|H|+|U|)\times d}$，再按来源拆成$\bm E_H\in\mathbb R^{|H|\times d}$与$\bm E_U\in\mathbb R^{|U|\times d}$。解码器每层在自注意力之后，对两部分分别做编码器--解码器注意力，得到$\bm c^H_t$与$\bm c^U_t$，拼接后送入前馈层。最后一层在第$t$步对历史词元和当前词元各给出一个注意力分布$\bm\alpha_t\in\mathbb R^{|H|}$与$\bm\beta_t\in\mathbb R^{|U|}$，输出分布为
\begin{equation}
  p(r_t=w\mid H,U,r_{<t})=\lambda_t\sum_{i:\,H_i=w}\alpha_{t,i}+(1-\lambda_t)\sum_{j:\,U_j=w}\beta_{t,j},
  \label{eq:nlp-analysis-two-source-copy}
\end{equation}
\begin{equation}
  \lambda_t=\sigma\bigl(\bm w_d^{\top}\bm d_t+\bm w_H^{\top}\bm c^H_t+\bm w_U^{\top}\bm c^U_t\bigr),
  \label{eq:nlp-analysis-copy-gate}
\end{equation}
其中$\bm d_t$是解码器当前状态，$\lambda_t\in(0,1)$是从历史复制的权重。$\lambda_t$接近0时模型照抄当前话语，接近1时从历史中取词；当前话语既无指代也无省略时，理想的行为是整句照抄、$\lambda_t$始终接近0。训练以参考改写的负对数似然为目标函数；这一模型无法生成历史与当前话语中都没有的词，因此它要求参考改写中的词都能在这两处找到。

用“那明天呢”在H1下的改写看一步计算（数值为教学构造）。解码到“明天能到”之后，下一个词应为“图书馆”。设$\lambda_t=0.8$，历史注意力在“图书馆”上为$0.9$，当前话语中没有这个词，则$p(\text{图书馆})=0.8\times0.9=0.72$；“明天”只在当前话语中出现，若当前注意力在它上面为$0.6$，则$p(\text{明天})=0.2\times0.6=0.12$。模型输出“图书馆”。若门控错误地给出$\lambda_t=0.3$，“图书馆”的概率降为$0.27$，与当前话语中某个词的概率相当，改写就可能停在“明天能到呢”这样的不完整句子上。门控因此承担了“此处是否需要从历史补内容”的判断，相当于隐式地定位了省略位置。

只做复制的设计限制了模型补入新内容，这既是优点也是限制：它天然满足“补全内容来自历史”这一要求，却无法生成连接词、语序调整所需的虚词。其他方法从不同角度处理同一问题。Pan等的方法分两步，先在历史中逐词标注哪些词应被补入，再把选出的词附在当前话语之后，交给一个指针生成网络写出完整句子，选词一步可以单独评价\scite{nlp-analysis-pan2019restoration}。选词写成历史词元上的二分类：用预训练编码器把历史与当前话语一起编码，第$i$个历史词元的表示为$\bm h_i$，补入概率为$\sigma(\bm w^{\top}\bm h_i)$；Liu等把改写看作在“当前话语词×历史词”的矩阵上预测编辑操作（在某处插入历史中的某个片段，或用历史片段替换当前话语中的某个词），用语义分割网络一次预测整个矩阵，避免了逐词解码的错误累积\scite{nlp-analysis-liu2020run}。与槽位继承相比，这些方法不需要槽位定义，覆盖面广，代价是输出不再是可以逐项检查的结构，补入的片段是否放在了正确的语法位置需要另外核对：“图书馆”被复制到人物的位置上，复制来源正确，改写仍然是错的。

\subsubsection{基于大语言模型的查询改写}
\label{sec:nlp-analysis-llm-rewriting}

前两种方法都需要成对的改写训练数据。大语言模型可以直接按提示完成改写：给出若干“历史、当前问题、独立问题”的示例，再给出当前对话，要求模型输出一个不依赖历史也能理解的问题。改写的要求可以写进提示：保留用户原意，补出省略和指代的内容，不添加历史中没有的信息，历史不足以补全时原样输出。

代码~\ref{code:nlp-analysis-rewrite-prompt}给出一个少样本改写提示的骨架，示例中同时包含一个无需改写和一个历史不足的情形，用来教模型在这两种情况下原样输出。

\begin{codeblock}[label={code:nlp-analysis-rewrite-prompt}]{少样本话语改写提示}
把“当前问题”改写为不看对话历史也能理解的独立问题。
只使用历史中出现过的信息；无需改写或历史不足时原样输出。

历史：今天能到图书馆借书吗？——可以，开到晚上九点。
当前问题：那明天呢？
独立问题：明天能到图书馆借书吗？

历史：（无）
当前问题：那明天呢？
独立问题：那明天呢？

历史：读书会改到周五下午两点了。
当前问题：读书会在哪里举行？
独立问题：读书会在哪里举行？

历史：{对话历史}
当前问题：{当前问题}
独立问题：
\end{codeblock}

在对话式检索中，改写的目的是检索效果，这带来两个新的做法。Mao等的LLM4CS让模型对同一问题生成多个改写，并生成一个假想的回答，把这些文本的嵌入加以聚合作为检索查询，以多个样本平均掉单次改写的偏差\scite{nlp-analysis-mao2023llm4cs}。Ye等提出“信息充分的改写”，要求改写不仅补全省略，还要补入历史中有助于检索的背景信息，并以正确、清晰、信息充分、不冗余四项要求编写提示；他们再用大模型的改写结果训练一个较小的改写模型，以降低推断成本\scite{nlp-analysis-ye2023informative}。

大语言模型改写的主要风险是补入历史中没有的内容。模型可能把“那明天呢”改写为“明天上午九点能到市图书馆借书吗”，“上午九点”“市”都是模型自己的猜测。一个简单而有效的检查是来源核对，见算法~\ref{alg:nlp-analysis-rewrite-check}：改写中每个不在当前话语中的实词，都应能在历史中找到对应；找不到来源的新增内容要么删去，要么交给人工或澄清流程。

\begin{algorithm}[改写结果的来源核对]
  \label{alg:nlp-analysis-rewrite-check}
  \begin{algorithmic}[1]
    \Require 对话历史$H$，当前话语$u$，候选改写$r$；可接受的规范化映射$N$（如“明天”到具体日期）
    \Ensure 通过核对的改写，或带标记的无来源片段
    \State 对$H$、$u$、$r$分词，去掉虚词与标点，得到实词集合$W_H$、$W_u$、$W_r$
    \State $D\gets W_r\setminus W_u$\Comment{改写中新增的实词}
    \For{每个$w\in D$}
      \If{$w\in W_H$，或$w$可由$N$从$W_H\cup W_u$中的某个词得到}
        \State 记录$w$的来源位置
      \Else
        \State 把$w$标记为无来源
      \EndIf
    \EndFor
    \If{没有无来源的词}
      \State \Return $r$
    \Else
      \State \Return 删去无来源片段后的$r$，或请求澄清，并附上被标记的词
    \EndIf
  \end{algorithmic}
\end{algorithm}

字符串层面的核对只能发现最明显的捏造：同义改述（历史说“借阅”，改写写“借书”）会被误报，需要配合同义词表或嵌入相似度；把正确来源的词放到错误位置，它也无法发现。本节关注的是语义补全，即改写是否忠实地还原了用户的问题；第~\ref{chap:nlp-question-answering-dialogue}章从检索与问答系统的角度讨论查询改写，那里的评价标准是检索效果，可能有意加入扩展词，两者的目标不同。

表~\ref{tab:nlp-analysis-rewriting-methods}对照三类方法。

\begin{table*}[htbp]
  \centering
  \small
  \begin{tabularx}{\linewidth}{@{}l XXX@{}}
    \toprule
    方法 & 解决的问题 & 核心机制 & 新的代价或局限\\
    \midrule
    槽位继承 & 结构化场景中的省略 & 逐槽判断历史槽值是否继承 & 只覆盖预定义槽位；槽间依赖需单独处理\\
    复制式改写 & 开放领域的任意省略 & 从历史与当前句两路复制，门控$\lambda$选择来源 & 需要改写训练数据；补入位置需另行核对\\
    大语言模型改写 & 无训练数据；面向检索的改写 & 少样本提示，多次采样聚合 & 可能补入无来源内容，需来源核对；推断成本高\\
    \bottomrule
  \end{tabularx}
  \caption{省略恢复与话语改写方法对照。}
  \label{tab:nlp-analysis-rewriting-methods}
\end{table*}

\subsection{会话含义与间接言语行为}
\label{sec:nlp-analysis-implicature}

上一小节补出的是说话者省略、但历史中明确出现过的内容。还有一类意义从未被说出，却是说话者想要传达的。问“同学们都到了吗”，回答“有些同学到了”，听话者通常会理解为“不是所有同学都到了”，尽管“有些”在逻辑上并不排除“全部”。问“你能把窗关上吗”，听话者知道这是请求关窗，而不是询问自己有没有关窗的能力。前者是会话含义，后者是间接言语行为，二者都要求在字面意义之上推断说话者的意图。

这类推断有两个挑战。其一，它依赖合作假设与背景：听话者之所以从“有些”推出“不是全部”，是因为假定对方提供了尽可能充分的信息；对方若不知道其余同学的情况，这一推断就不成立。其二，推断结果本身可能有分歧，不同标注者对同一句话是否带有某种含义的判断并不总是一致。方法的演进从语言学的规则与分类模型，过渡到用自然语言推断的形式检验模型的判断，再到对大语言模型的语用推理能力做专门评测。评测时的核心问题是区分模型是在推理含义，还是在匹配训练数据中的表面模式。

\subsubsection{合作原则与含义推理}
\label{sec:nlp-analysis-grice}

Grice提出，交谈的双方默认遵守\term{合作原则}（cooperative principle）：使自己的话语在当时的交谈目的下恰如其分\scite{nlp-analysis-grice1975}。合作原则展开为四条准则。量的准则：提供当前交谈所需的信息，不多也不少。质的准则：不说自己认为不真的话，不说缺乏证据的话。关系准则：说话要相关。方式准则：避免晦涩和歧义，简洁而有条理。这些准则描述的不是人们实际总是怎样说话，而是交谈双方默认对方会遵守的期望，听话者据此解释话语。当字面意义看起来违反某条准则，而听话者仍相信对方在合作时，就会推出一个使话语恢复合作的额外意义，这就是\term{会话含义}（conversational implicature）。

等级含义是最常见、也最常被计算研究的一类。“有些”与“全部”构成一个按信息强度排序的等级$\langle\text{有些},\text{全部}\rangle$，“全部”蕴含“有些”而反之不然。推理分四步：说话者说了较弱的“有些同学到了”；若全部到了，说较强的“全部到了”信息更充分，按量的准则说话者应当这样说；说话者没有这样说，又假定他在合作且了解情况；因此他认为“全部到了”不成立。“或”与“和”、“可能”与“一定”、“温”与“热”构成类似的等级。关系含义则依赖关系准则：甲问“还有咖啡吗”，乙答“楼下有便利店”，字面上答非所问，听话者假定它相关，推出“这里没有咖啡了，可以去楼下买”。

两类含义的依赖程度不同。Grice把不依赖特定语境、只要用了某种说法就默认出现的含义称为\term{一般会话含义}（generalized conversational implicature），把只在特定语境中才出现的含义称为\term{特殊会话含义}（particularized conversational implicature）。等级含义通常归入前者：只要说了“有些”，在没有相反信息时就推出“不是全部”。关系含义属于后者：“楼下有便利店”在另一场对话中不会有“没有咖啡了”的意思。前者有规律可循，适合用模板构造评测数据，后者必须结合具体对话判断。说话者还可能故意明显地违反某条准则，让听话者意识到违反并推出另一层意义：天下着大雨，有人说“真是出门的好天气”，明显违反质的准则，听话者据此推出反讽；推荐信只写“该生出勤良好、字迹工整”，明显违反量的准则，读者推出“学术能力不值一提”。反讽与言外之意的识别因此也是会话含义推理的一部分。

等级含义的推理可以写成一个可计算的概率模型。理性言语行为（rational speech act, RSA）模型把听话者与说话者的相互推理写成递归\scite{nlp-analysis-frank-goodman2012}。设可能的世界状态为$w$，话语为$u$，$[\![u]\!](w)\in\{0,1\}$表示$u$在$w$中字面上是否为真，$P(w)$为先验。字面听话者只按真值过滤：$L_0(w\mid u)\propto[\![u]\!](w)P(w)$。说话者按听话者能否猜对来选择话语：$S_1(u\mid w)\propto\exp\bigl(\alpha\log L_0(w\mid u)\bigr)$，$\alpha>0$衡量说话者的理性程度。语用听话者再推断说话者为什么选了这句话：$L_1(w\mid u)\propto S_1(u\mid w)P(w)$。

取两个世界“部分到了”与“全部到了”，先验各$0.5$，两句话“有些到了”（两个世界中都为真）与“全部到了”（只在后者中为真），$\alpha=1$。字面听话者听到“有些”时两个世界各$0.5$，听到“全部”时只剩“全部到了”。在“全部到了”的世界里，说话者说“有些”能让听话者猜对的概率为$0.5$，说“全部”为1，于是$S_1(\text{有些}\mid\text{全部})=0.5/1.5=1/3$；在“部分到了”的世界里，“全部”为假，$S_1(\text{有些}\mid\text{部分})=1$。语用听话者听到“有些”后，$L_1(\text{部分}\mid\text{有些})\propto1\times0.5$，$L_1(\text{全部}\mid\text{有些})\propto\frac13\times0.5$，归一化得$0.75$与$0.25$。字面意义对两个世界不作区分，递归一次之后，“不是全部”成为更可能的解读，这正是上面四步推理的概率版本；$\alpha$越大，偏向越明显。RSA的局限在于需要事先列出世界状态与备选话语，适合控制良好的语言游戏，难以直接用于开放文本。

会话含义与逻辑蕴含的区别在于\term{可取消性}（cancellability）：含义可以被后续话语明确否定而不产生矛盾，蕴含不能。“有些同学到了，实际上全部都到了”虽然略显迂回，但并不自相矛盾，说明“不是全部”只是含义；“全部同学都到了，实际上有些没到”则自相矛盾，说明“全部”蕴含“有些”。这一检验给出了标注与评测时区分两者的可操作标准。

\term{预设}（presupposition）是第三类需要区分的意义：它是话语成立所依赖的背景前提。“小王停止抽烟了”预设小王曾经抽烟；“小王知道会议改期了”预设会议确实改期了。引出预设的词语称为预设触发语，常见的有：有定描述（“小王的车”预设小王有车）、叙实动词（“知道”“意识到”）、状态改变动词（“不再”“开始”“停止”）、反复类副词（“又”“再”）以及分裂句（“是小王打碎了杯子”预设有人打碎了杯子）。预设的特征是在否定、疑问、条件等算子下仍然保留：“小王没有停止抽烟”“小王停止抽烟了吗”“如果小王停止抽烟”都仍然预设他曾经抽烟，而蕴含在否定下不保留。预设也不是完全不可取消：在否定句中，说话者可以明确否认它，“小王没有停止抽烟，他根本就没抽过烟”并不矛盾；但在肯定句中加上同样的否认就会自相矛盾。表~\ref{tab:nlp-analysis-meaning-types}用两条检验区分三种意义。

\begin{table}[htbp]
  \centering
  \small
  \begin{tabularx}{\linewidth}{@{}l X cc@{}}
    \toprule
    意义类型 & 例子 & 否定后是否保留 & 能否取消\\
    \midrule
    蕴含 & “全部同学到了”$\Rightarrow$“有些同学到了” & 否 & 否\\
    会话含义 & “有些同学到了”$\leadsto$“不是全部” & 否 & 能\\
    预设 & “小王停止抽烟”$\gg$“小王曾经抽烟” & 是 & 肯定句中不能\\
    \bottomrule
  \end{tabularx}
  \caption{蕴含、会话含义与预设的区分检验。$\leadsto$表示含义，$\gg$表示预设。}
  \label{tab:nlp-analysis-meaning-types}
\end{table}

\subsubsection{间接言语行为识别}
\label{sec:nlp-analysis-indirect-acts}

Austin指出，说话本身就是在做事：说“我保证明天交稿”就完成了一次许诺，这样的话语没有真假，只有是否得当\scite{nlp-analysis-austin1962}。Searle把这一思想系统化为言语行为理论，并用一组\term{适切条件}（felicity conditions）刻画每种行为成功实施的条件。以请求听话者H做动作A为例，四类条件为：命题内容条件，A是H将来的动作；准备条件，H有能力做A，且不请求的话H不一定会去做；真诚条件，说话者想要H做A；本质条件，这句话算作让H做A的一次尝试。Searle后来又按行为的目的把言语行为分为五类：断言类（陈述、报告，使话语与世界相符）、指令类（请求、命令、提问，让听话者去做某事）、承诺类（许诺、答应，说话者约束自己将来的行为）、表达类（感谢、道歉，表达心理状态）、宣告类（任命、宣布，话语本身改变了世界的状态）。第~\ref{sec:nlp-analysis-dialogue-acts}小节的话语行为标签集，大体可以看作这一分类在对话中的细化，例如请求与是非问都属于指令类，感谢与附和更接近表达类。

\term{间接言语行为}（indirect speech act）是通过实施一种行为来完成另一种行为。Searle观察到，间接请求的句式系统地对应于对适切条件的询问或陈述\scite{nlp-analysis-searle1975}。“你能把窗关上吗”询问的是准备条件中H的能力；“我希望你把窗关上”陈述的是真诚条件；“你会把窗关上吗”询问的是命题内容条件中H将来的动作。听话者的推理过程大致是：字面上这是一个关于能力的问题；在当前情境中，对方显然知道我有能力关窗，这个问题作为信息询问没有意义；询问能力是请求的准备条件之一；因此对方很可能是在请求我关窗。这里同样依赖合作假设，并且推理的关键一步是“字面解读在情境中不合理”。若情境改变，例如在讨论受伤后的康复情况，“你能把窗关上吗”就可能是真的在问能力。

把这一推理拆成可检查的步骤，便于看出它在哪里依赖情境：
\begin{enumerate}
  \item 确定字面行为：话语的句式是疑问，字面上询问H是否有能力做A。
  \item 检查字面解读是否合作：在当前情境中，S是否已经知道答案；若已知，作为信息询问就违反了量或关系的准则。
  \item 寻找被询问或陈述的适切条件：能力属于请求的准备条件。
  \item 推断间接行为：S通过询问准备条件来实施请求。
  \item 检查情境是否支持该请求：A是否对S有益、H是否方便做A；不支持时退回字面解读。
\end{enumerate}
第2步与第5步都需要情境知识，这也是计算方法最难自动获得的部分。

计算上的识别方法经历了两个阶段。Perrault与Allen用基于规划的方法建模这一推理：把言语行为表示为带前提与效果的规划算子，听话者从字面行为出发，推断说话者可能的目标，再识别出间接的请求\scite{nlp-analysis-perrault-allen1980}。这种方法可解释，但需要手工编写情境知识。统计方法把识别写成条件分类：输入当前话语和可见的上下文，输出它实际承担的行为，即$p(a\mid u,\text{上下文})$，与第~\ref{sec:nlp-analysis-dialogue-acts}小节的话语行为分类同形，区别只在于标签表示的是实际意图而不是句式。间接回答是一个重要的特例：问“你周末要来聚会吗”，回答“我得赶一篇论文”，要识别的是“不来”。识别时若上下文不足以排除字面解读，模型应当保留字面解读，而不是默认存在间接意图；“这里有点冷”在讨论天气时只是陈述感受，在对方站在窗边时才可能是请求关窗。

\subsubsection{语用推理能力的评测}
\label{sec:nlp-analysis-pragmatic-eval}

语用推理很难直接标注，评测通常把它转写为已有任务的形式。IMPPRES把含义与预设写成自然语言推断问题\scite{nlp-analysis-jeretic2020imppres}。对等级含义，给出前提“有些同学到了”与假设“不是所有同学都到了”，按语用解读应判为蕴含，按逻辑解读应判为中立，模型的选择反映它采用哪一种解读。对预设，数据把触发语放在否定、疑问、情态、条件等算子下，检验模型是否仍然推出预设。数据由模板生成，规模可控，便于逐现象分析。Jeretic等发现，在MultiNLI上训练的BERT稳定地把“some”引出的等级含义当作蕴含，即采用语用解读；对部分预设触发语（如only），即使触发语处在否定之下，它也能推出预设，对另一些触发语则不能。

Circa检验间接回答的理解\scite{nlp-analysis-louis2020circa}。数据含约3.4万个是非问句与间接回答对，由众包人员在给定情境下写出问题和不直接说“是”或“否”的回答，标签包括是、否、有条件的是、介于两者之间等。模型需要从回答推断它对问题的立场。Louis等的基线把问题与回答拼接后送入BERT做句对分类，并比较了先在自然语言推断数据MultiNLI上微调、再在Circa上微调的做法；后者的准确率最高，但与不经迁移的模型相比差异在统计上不显著，因此不能据此断言推断数据带来了稳定收益。他们还比较了只给回答、不给问题的设置：准确率显著下降，说明模型确实利用了问题，但下降后的准确率仍然不低，相当一部分样本单凭回答的表面形式就能判断。在训练中未出现过的情境上测试时，准确率也会下降。

对大语言模型，Ruis等设计了二选一的会话含义评测：把一问一答写成“甲问……，乙答……，乙的意思是”，比较模型续写“是”与“否”的概率\scite{nlp-analysis-ruis2023}。例如“甲问：你周末要来聚会吗？乙答：我得赶一篇论文。乙的意思是”，若模型给“否”的概率为$0.7$、“是”为$0.3$，记为答对。同一样本换用措辞不同的六个模板，若三个模板答对、三个答错，单一模板的结论就完全取决于选了哪一个，因此应报告在所有模板上的平均准确率及其范围。他们比较了只做预训练、对话微调、基准级指令微调与样例级指令微调四类模型，结论是：只做预训练的模型在零样本下接近随机水平，样例级指令微调的模型明显更好，少样本示例进一步提高准确率，但多数模型仍低于人类的一致率。在他们的设置中，六个模板之间的波动不大；这一稳健性依赖于对多个模板取平均的评测设计，换用新的格式时仍需重新检验。Hu等用多项选择形式比较人与语言模型在讽刺、隐喻、间接言语、违反准则等七类现象上的表现，并分析模型的错误类型与人是否相近\scite{nlp-analysis-hu2023pragmatics}。

这些评测的设计给出几条经验。其一，结果可能随提示措辞和示例变化，报告时应给出多个模板的均值与波动，而不是只报最好的模板。其二，需要人类基线：人对同一问题的判断本身存在分歧，模型的准确率应与人类的一致程度对照看。其三，要设计对照样本以区分推理与模式匹配：保持回答不变、改变问题或情境，若模型的判断不随之改变，说明它依赖的是回答的表面形式，而不是推理。语用评测的分数衡量的是模型在特定格式下给出符合多数人解读的答案的能力，能否推广到开放对话中的语用理解，需要结合具体应用检验。
